
Évolution d'un système physique
Opérateur d'évolution
Comment  distinguer les états physiques ?
-États
-États
Problématique générale de la mesure :
-interaction système/instrument
-interprétation de Copenhague, interprétation d'Everett

%II.3
Dynamique d'un système

\subsection{Opérateur d'évolution}
\begin{center}
L'état à l'instant t $\to$ état à l'instant t'
\end{center}
\[
|\Psi(t)>\qquad\xrightarrow[U(t,t')]{}\quad|\Psi(t')>
\]
\subsection{Propriétés générale}
\subsubsection{linéarité}
\[
|\Psi_1(t)>\quad\to\quad|\Psi_1(t')> = U(t,t')|\Psi_1(t)>
\]
\[
|\Psi_2(t)>\quad\to\quad|\Psi_2(t')> = U(t,t')|\Psi_2(t)>
\]
\[
|\Phi(t)>=\alpha_1|\Psi_1(t)>+\alpha_2|\Psi_2(t)>\quad\to\quad
|\Phi(t')>=\alpha_1|\Psi_1(t')>+\alpha_2|\Psi_2(t')> = U(t,t')|\Phi(t)>
\]

\subsubsection{Invariance par translation dans le temps}

\[
U(t,t') \to U(t'-t)
\]

\subsubsection{Continuité}
Si on part de deux états "voisins", on arrive à deux états voisins.

\subsubsection{Unitarité}

\section{Comment distinguer les états et les identifier}

\subsection{Via la valeur d'une grandeur physique}

\subsection{Problématique nouvelle en physique quantique}

système $\mc{S}$, grandeur G 

